\begin{derivation}[从\eqnrefs{eq:qcd-singlet-axial-rotation-dp68}到\eqnrefs{eq:qcd-singlet-axial-anomaly-dp32}]
在欧几里得时空中取规范表示 $R$，并令
\begin{equation*}
 D_\mu=\partial_\mu+\ii g\mathcal A_\mu^aT^a,
 \qquad
 \Tr_R(T^aT^b)=T(R)\delta^{ab}.
\end{equation*}
用厄米算符 $\mathcal H\equiv\ii\slashed D_E$ 的正交本征函数展开 Grassmann 场：
\begin{equation*}
 \mathcal H\phi_n=\lambda_n\phi_n,
 \qquad
 \psi=\sum_na_n\phi_n,
 \qquad
 \bar\psi=\sum_n\bar b_n\phi_n^\dagger.
\end{equation*}
为了定义量子测度沿轴向方向的变化，取一参数局域轴向轨道
\begin{equation*}
 \psi_\tau=\ee^{\ii\tau\alpha(x)\gamma^5}\psi,
 \qquad
 \bar\psi_\tau=\bar\psi\,\ee^{\ii\tau\alpha(x)\gamma^5},
 \qquad \tau\in\mathbb R.
\end{equation*}
在本征函数系数空间中，$\psi$ 的变换矩阵为
\begin{equation*}
 C_{mn}(\tau)
 \equiv\int\dd^4x\,
 \phi_m^\dagger(x)\ee^{\ii\tau\alpha(x)\gamma^5}\phi_n(x),
 \qquad C(0)=I,
\end{equation*}
其在恒等元处的切向量严格为
\begin{equation*}
 \left.\frac{\dd C_{mn}}{\dd\tau}\right|_{\tau=0}
 =\ii\int\dd^4x\,
 \alpha(x)\phi_m^\dagger\gamma^5\phi_n.
\end{equation*}
Grassmann 测度按行列式的逆变换；$\psi$ 与 $\bar\psi$ 两份测度相乘。令 $J(\tau)$ 为总 Jacobian，则
\begin{align*}
 \left.\frac{\dd}{\dd\tau}\ln J(\tau)\right|_{\tau=0}
 &=-2\Tr\!\left[\left.\frac{\dd C}{\dd\tau}\right|_{0}\right]\\
 &=-2\ii\int\dd^4x\,\alpha(x)
 \sum_n\phi_n^\dagger(x)\gamma^5\phi_n(x).
\end{align*}
Ward 恒等式只需要这一沿群轨道的切向变化；随后对迹作规范协变正则化。
模和紫外发散，必须用保持规范协变性的同一算符谱正则化：
\begin{equation*}
 \mathscr D_R(x)
 \equiv\lim_{M\to\infty}
 \tr_{\rm spin,R}\!\left[
 \gamma^5\langle x|\ee^{-\mathcal H^2/M^2}|x\rangle
 \right].
\end{equation*}
采用欧氏 $\gamma$ 矩阵与曲率约定，Dirac 算符平方写成
\begin{equation*}
 \mathcal H^2=-D^2+\frac g2\sigma^{\mu\nu}\mathcal F_{\mu\nu}^aT^a.
\end{equation*}
在固定点 $x$ 插入平面波完备集，把 $D_\mu$ 平移成 $D_\mu+\ii p_\mu$。$M\to\infty$ 时只需保留能在高斯积分后留下 $M^0$ 的项。零阶项含 $\tr\gamma^5=0$；一次曲率项含 $\tr(\gamma^5\sigma^{\mu\nu})=0$；第一项非零贡献来自二次曲率：
\begin{align*}
 \mathscr D_R(x)
 &=\lim_{M\to\infty}\Bigg\{
 \int\frac{\dd^4p}{(2\pi)^4}\ee^{-p^2/M^2}
 \frac12\left(\frac{g}{2M^2}\right)^2
 \tr_{\rm spin}\!\left(\gamma^5\sigma^{\mu\nu}\sigma^{\rho\sigma}\right)
 \Tr_R(T^aT^b)
 \mathcal F_{\mu\nu}^a\mathcal F_{\rho\sigma}^b
 +\mathscr R_M(x)\Bigg\},
\end{align*}
其中 $\mathscr R_M(x)$ 定义为减去所显示二次曲率项后剩余的全部热核贡献，并满足
\begin{equation*}
 \lim_{M\to\infty}\mathscr R_M(x)=0.
\end{equation*}
使用
\begin{equation*}
 \int\frac{\dd^4p}{(2\pi)^4}\ee^{-p^2/M^2}=\frac{M^4}{16\pi^2},
 \qquad
 \tr_{\rm spin}(\gamma^5\sigma^{\mu\nu}\sigma^{\rho\sigma})
 =-4\epsilon^{\mu\nu\rho\sigma},
\end{equation*}
以及 ${}^\star\!\mathcal F^{\mu\nu}\equiv\epsilon^{\mu\nu\rho\sigma}\mathcal F_{\rho\sigma}/2$，所有 $M$ 次幂恰好抵消。按正文采用的欧氏到闵氏延拓与 $\epsilon^{0123}=+1$ 约定，得到
\begin{equation*}
 \boxed{
 \mathscr D_R(x)
 =\frac{g^2T(R)}{16\pi^2}
 \mathcal F_{\mu\nu}^a{}^\star\!\mathcal F^{a\mu\nu}.}
\end{equation*}
把这个极限代回测度的切向变化
\begin{equation*}
 \left.\frac{\dd}{\dd\tau}\ln J(\tau)\right|_{0}
 =-2\ii\int\dd^4x\,\alpha(x)\mathscr D_R(x),
\end{equation*}
再与质量项沿同一轴向轨道的经典变化相加，Ward 恒等式为
\begin{equation*}
 \partial_\mu J_5^\mu
 =\frac{2\ii mc}{\hbar}\bar\psi\gamma^5\psi
 +\frac{g^2\,2T(R)}{16\pi^2}
 \mathcal F_{\mu\nu}^a{}^\star\!\mathcal F^{a\mu\nu}.
\end{equation*}
QCD 基本表示满足 $T_F=1/2$；对 $N_f$ 个 Dirac 味求和即得到正文\eqnrefs{eq:qcd-singlet-axial-anomaly-dp32}。反常系数由热核在 $M\to\infty$ 后留下的有限曲率二次项固定；规范协变调节函数之间的差异进入 $\mathscr R_M(x)$，而该余项在去除调节器的极限中消失。
\end{derivation}

\begin{derivation}[从\eqnrefs{eq:electroweak-topological-charge-dp68}到\eqnrefs{eq:sm-bplusl-selection-dp32}]
先把拓扑信息与费米子路径积分连接起来。对一个处在 $SU(2)$ 表示 $R$ 的左手 Weyl 场，欧几里得 Dirac 算符的指数为
\begin{equation*}
 \operatorname{ind}\slashed D_E
 \equiv n_L-n_R=2T(R)\nu.
\end{equation*}
标准模型弱二重态属于 $R=\mathbf2$，$T(\mathbf2)=1/2$。当 $\nu>0$ 时，每个左手二重态有 $\nu$ 个左手零模而没有相应右手零模：
\begin{equation*}
 n_L-n_R=\nu.
\end{equation*}

把一个左手二重态沿 Dirac 本征函数展开，显式分离零模系数，
\begin{equation*}
 \psi_L(x)=\sum_{r=1}^{\nu}a_r\phi_{0r}(x)
 +\sum_{\lambda_n\neq0}a_n\phi_n(x).
\end{equation*}
费米子二次作用量只含非零本征值，因而测度中零模方向留下独立 Grassmann 积分
\begin{equation*}
 \int\mathcal D\psi_L\,\ee^{-S_F}
 =\mathcal N_{\rm nzm}
 \left(\prod_{r=1}^{\nu}\int\dd a_r\right),
 \qquad
 \mathcal N_{\rm nzm}\equiv
 \mathcal N_{\rm meas}\prod_{\lambda_n\neq0}\lambda_n,
\end{equation*}
Grassmann 规则 $\int\dd a_r=0$、$\int\dd a_r\,a_r=1$ 表明：没有外场插入时该拓扑扇区的费米子积分为零；一个非零相关函数必须为每个零模恰好提供一个相应的场因子。对 $\nu=1$，每个左手弱二重态必须出现一次。

一代标准模型含三个颜色副本 $Q_L^a$（$a=1,2,3$）和一个 $L_L$。把颜色收缩成单态、弱同位旋指标用 $\epsilon_{ij}$ 配对，零模饱和产生的局域多费米子结构可示意写成
\begin{equation*}
 \mathfrak O_{\rm tH}^{(g)}
 \equiv
 \epsilon_{abc}
 \bigl(Q_{Lg}^{a}Q_{Lg}^{b}\bigr)
 \bigl(Q_{Lg}^{c}L_{Lg}\bigr),
\end{equation*}
其中括号表示 Lorentz 与 $SU(2)_L$ 指标按不变张量收缩。三个夸克场各带 $B=1/3$，一个轻子场带 $L=1$，所以单代、单位拓扑荷的顶点改变
\begin{equation*}
 \Delta B_{g}=3\times\frac13=1,
 \qquad
 \Delta L_{g}=1.
\end{equation*}

同一结果也可直接从反常 Ward 恒等式积分得到。对重子流和轻子流，零模计数给每代相同的拓扑密度系数，
\begin{equation*}
 \partial_\mu j_B^\mu
 =N_g\frac{g^2}{32\pi^2}
 \mathcal W_{\mu\nu}^I{}^\star\!\mathcal W^{I\mu\nu},
 \qquad
 \partial_\mu j_L^\mu
 =N_g\frac{g^2}{32\pi^2}
 \mathcal W_{\mu\nu}^I{}^\star\!\mathcal W^{I\mu\nu}.
\end{equation*}
在一个连接两个渐近真空的四维区域上积分，并使用正文对 $\nu$ 的定义，
\begin{align*}
 \Delta B
 &=\int\dd^4x\,\partial_\mu j_B^\mu
 =N_g\nu,\\
 \Delta L
 &=\int\dd^4x\,\partial_\mu j_L^\mu
 =N_g\nu.
\end{align*}
对任意整数 $\nu$，$|\nu|$ 个零模副本给出同样的线性倍数；$\nu<0$ 时手征性与粒子数变化同时反向。于是
\begin{equation*}
 \Delta(B-L)=0,
 \qquad
 \Delta(B+L)=2N_g\nu,
\end{equation*}
即正文\eqnrefs{eq:sm-bplusl-selection-dp32}。选择定则来自费米子测度在不同拓扑扇区中的零模结构，而不是规范反常未抵消。
\end{derivation}

\begin{derivation}[从\eqnrefs{eq:qcd-singlet-axial-anomaly-dp32}到\eqnrefs{eq:strong-cp-theta-bar-dp11}]
把所有有色 Dirac 夸克收成味向量 $q=(u,d,s,\ldots)^T$，质量项写成
\begin{equation*}
 \Lag_m=-\bar q_LMq_R-\bar q_RM^\dagger q_L,
\end{equation*}
其中 $M$ 可以是任意非奇异复矩阵。取对角的实味空间矩阵
\begin{equation*}
 A={\rm diag}(\alpha_1,\ldots,\alpha_{N_f})
\end{equation*}
并作独立轴向重相位
\begin{equation*}
 q_L\longrightarrow \ee^{-\ii A}q_L,
 \qquad
 q_R\longrightarrow \ee^{+\ii A}q_R.
\end{equation*}
因为 $\bar q_L\to\bar q_Le^{+\ii A}$，质量矩阵在新场基底中变成
\begin{equation*}
 M' =\ee^{+\ii A}Me^{+\ii A}.
\end{equation*}
对行列式取相位，利用 $\det(\ee^{\ii A})=\ee^{\ii\Tr A}$：
\begin{align*}
 \det M'
 &=\ee^{\ii\Tr A}\det M\,\ee^{\ii\Tr A}
 =\ee^{2\ii\Tr A}\det M,\\
 \arg\det M'
 &=\arg\det M+2\Tr A
 \pmod{2\pi}.
\end{align*}
因此轴向旋转确实可以移动夸克质量矩阵的整体复相位；矢量型重相位 $q_L,q_R\to \ee^{\ii B}q_{L,R}$ 则不会产生这一相位变化。

另一方面，正文\eqnrefs{eq:qcd-singlet-axial-anomaly-dp32} 表明同一个轴向变换不是量子测度的对称。正文\eqnrefs{eq:qcd-theta-action-dp89} 已固定 $\theta_s$ 的归一化。对第 $i$ 个基本表示 Dirac 味作角度 $\alpha_i$ 的轴向旋转，Fujikawa Jacobian 给
\begin{equation*}
 \ln J_i
 =-\ii\alpha_i\frac{g_s^2}{16\pi^2}
 \int\dd^4x\,
 \mathcal F_{\mu\nu}^a{}^\star\!\mathcal F^{a\mu\nu}.
\end{equation*}
把 $J_i$ 吸收到 $\exp(\ii S/\hbar)$ 中，等价于
\begin{equation*}
 \theta_s\longrightarrow\theta_s-2\alpha_i.
\end{equation*}
全部味同时变换时便有
\begin{equation*}
 \theta_s' =\theta_s-2\Tr A.
\end{equation*}
与质量行列式的变化相加：
\begin{align*}
 \theta_s'+\arg\det M'
 &=(\theta_s-2\Tr A)
 +(\arg\det M+2\Tr A)\\
 &=\theta_s+\arg\det M.
\end{align*}
所以轴向基底选择能分别改变 $\theta_s$ 与质量矩阵的整体相位，却不能改变组合
\begin{equation*}
 \boxed{\bar\theta=\theta_s+\arg\det M}.
\end{equation*}
对标准模型六个夸克，把 $M$ 分成上、下型两个块就得到正文\eqnrefs{eq:strong-cp-theta-bar-dp11} 的 $\arg\det(M_uM_d)$。

若某个夸克严格无质量，例如某个质量本征值 $m_j=0$，则相应轴向角 $\alpha_j$ 不会在质量项中产生可观测相位：$\det M=0$ 时整体行列式相位本身不再是物理参数。此时可以选择 $\alpha_j=\theta_s/2$ 把拓扑角完全移入这个无质量场的轴向定义，从而使 $\theta_s$ 不可观测。真实 QCD 中所有夸克质量均非零，因此这一消除机制不存在，$\bar\theta$ 成为独立的强 $CP$ 参数。
\end{derivation}
